\documentclass[conference]{template/IEEEtran}

\usepackage{amsmath}
\usepackage{amssymb}
\usepackage{array}
\usepackage{url}

\title{A Port-Aware and Concurrent Path Generator for Conditional Data Flow Diagrams}

\author{
\IEEEauthorblockN{Author Name}
\IEEEauthorblockA{
Course Project\\
School or Department\\
Email: author@example.com}
}

\begin{document}

\maketitle

\begin{abstract}
Conditional Data Flow Diagrams (CDFDs) describe software behavior through data flows, processes, control conditions, and hierarchical process decomposition. Generating paths from a CDFD is therefore not equivalent to enumerating ordinary graph paths. A process may have several input and output ports, different ports may be mutually exclusive, inputs within one port may need to be available together, and independent branches may run concurrently. This paper presents a CDFD path generation tool that accepts a complete CDFD project file, converts it into a common graph model, checks structural validity, generates atomic and concurrent paths, and visualizes both the graph and the resulting path information. The paper defines the path model, explains JSON and SOFL `.cdfd' parsing, describes legality checks and system functions, details the port-aware algorithms for different CDFD cases, and evaluates the tool on representative CDFD inputs.
\end{abstract}

\begin{IEEEkeywords}
CDFD, SOFL, path generation, concurrent path, graph analysis, port semantics
\end{IEEEkeywords}

\section{Background and Significance}

Conditional Data Flow Diagrams (CDFDs) are used in SOFL-style specifications~\cite{ipa_sofl_report} to describe how data items enter a system, are transformed by processes, are stored, and finally leave the system. Compared with a plain directed graph, a CDFD has richer semantics. Edges may represent data flow, active data flow, or control flow. Nodes may be processes, data stores, state nodes, condition nodes, or SOFL connector structures. A high-level process may also be decomposed into a lower-level CDFD.

The course task is to design a tool that can import an arbitrary CDFD following the agreed file format and automatically output its paths. The task requires the input file to contain all information needed by the tool, so the result should not depend on manual interpretation after import. It also requires a clear path definition because CDFD paths may involve concurrency and mutually exclusive process interfaces.

This topic is significant for two reasons. First, it turns an informal diagram-reading activity into a defined graph-analysis problem. A reader can ask whether a generated path is legal, why two branches are parallel, or why one input is insufficient to activate a process, and the tool can answer through explicit model rules. Second, it supports CDFD review and teaching. Instead of manually following arrows in a complex diagram, students can compare the imported CDFD, the generated atomic paths, concurrent paths, path relations, and consistency warnings.

\section{Method Overview}

This section defines the basic objects used by the tool and explains the overall method: input design, CDFD parsing, validity checking, path generation, and system output.

\subsection{CDFD Project Definition}

\textbf{Definition 1 (CDFD project).}
A CDFD project is a tuple
\begin{equation}
  P = (M, \mathcal{S}, \mathcal{G}, g_0),
\end{equation}
where \(M\) is the module information, \(\mathcal{S}\) is the set of process specifications, \(\mathcal{G}\) is the set of named CDFD graphs, and \(g_0 \in \mathcal{G}\) is the entry graph.

\textbf{Definition 2 (CDFD graph).}
A CDFD graph is a tuple
\begin{equation}
  G = (V, E, V_s, V_t, R),
\end{equation}
where \(V\) is the node set, \(E\) is the directed edge set, \(V_s\) is the set of source nodes, \(V_t\) is the set of sink nodes, and \(R\) is the set of explicit structures such as parallel, fork, choice, join, merge, and process decomposition hints.

\textbf{Definition 3 (edge).}
An edge \(e \in E\) has a source \(src(e)\), a target \(dst(e)\), a kind \(kind(e)\), a data label set \(data(e)\), and an optional condition \(cond(e)\). Data-flow and active-flow edges are traversable by paths. Control-flow edges are not path segments; they are collected as path conditions for the affected process.

\subsection{Input File Design}

The canonical project format is `cdfd-json-v1'. It contains:

\begin{itemize}
\item `module': constants, types, variables, and entry behavior graph;
\item `processes': process specifications, including input/output data, ports, preconditions, postconditions, and optional decomposition graph names;
\item `graphs': graph layers, each with nodes, edges, starts, ends, and explicit structures;
\item `metadata': optional layout or source-specific information.
\end{itemize}

The tool also imports SOFL desktop `.cdfd' XML files. A `.cdfd' file stores a single drawing as `componentList' and `connectionList'. The importer maps these XML elements into the same internal CDFD project model used by JSON. Thus the algorithm works on one common model rather than on several unrelated input formats.

YAML and CSV are not accepted as CDFD input formats. YAML would only provide another syntax for the same JSON contract. CSV cannot completely represent modules, process ports, decomposition, control conditions, and CDFD structures.

\subsection{Path Definitions}

\textbf{Definition 4 (process port).}
For a process \(p\), an input port \(q^{in}\) is a group of input data items or input edges. An output port \(q^{out}\) is a group of output data items or output edges. Inputs inside one input port use AND semantics. Different input ports are alternatives unless the CDFD explicitly states otherwise. Different output ports are also alternatives by default, while multiple edges in one selected output port may be produced together.

\textbf{Definition 5 (port readiness).}
Let \(D\) be the set of available data items and \(A\) the set of activated upstream nodes. An input port \(q\) of process \(p\) is ready if every non-control input edge assigned to \(q\) has its data available and its source node activated:
\begin{equation}
  Ready(p,q,D,A) \Leftrightarrow
  \forall e \in InEdges(p,q):
  data(e) \subseteq D \land src(e) \in A.
\end{equation}
If a port is declared with `mode = any', one ready input item is sufficient. The default mode is `all'.

\textbf{Definition 6 (process activation).}
A process \(p\) is activatable when at least one of its input ports is ready:
\begin{equation}
  Fire(p,D,A) \Leftrightarrow
  \exists q \in \Pi^{in}_p: Ready(p,q,D,A).
\end{equation}
If no explicit ports are defined, the process input list is treated as one AND port.

\textbf{Definition 7 (atomic path).}
An atomic path is a source-to-sink data-flow trace
\begin{equation}
  \pi = \langle v_0, e_1, v_1, \ldots, e_n, v_n \rangle,
\end{equation}
where \(v_0 \in V_s\), \(v_n \in V_t\), \(src(e_i)=v_{i-1}\), \(dst(e_i)=v_i\), each \(e_i\) is traversable, and every process activation in the trace satisfies Definition 6.

\textbf{Definition 8 (concurrent path).}
A concurrent path is a structured path term
\begin{equation}
  C ::= v \mid C_1 ; C_2 \mid Par(C_1,\ldots,C_k) \mid Xor(C_1,\ldots,C_k).
\end{equation}
Sequential composition is printed as `->', parallel composition as `[A || B]', and exclusive alternatives through path relations as `XOR'. A legal concurrent path must have legal atomic path projections and must preserve fork, join, and port-activation constraints.

\subsection{CDFD Parsing}

The JSON parser validates the file against the JSON Schema~\cite{json_schema} and then creates `CDFDProject', `CDFDGraph', `Node', `Edge', `PortSpec', and `GraphStructure' objects. The SOFL parser reads XML components, resolves endpoints by component names and `shapeIndex', falls back to coordinate-based endpoint inference when a visible line touches a component, and preserves layout, port counts, connector indexes, and edge kinds as metadata.

After parsing, both JSON and SOFL inputs are represented as the same graph model:
\begin{equation}
  Input \rightarrow CDFDProject \rightarrow Path\ Algorithms.
\end{equation}

\subsection{Legality and Consistency Checks}

The tool performs two levels of validation. Schema validation is a hard check: if the JSON structure does not match the agreed format, the file is rejected. CDFD consistency checks are warnings because early CDFD models may be incomplete. Current checks include:

\begin{itemize}
\item process nodes without process specifications;
\item process specifications unused by any graph node;
\item graph data flows not declared in module variables;
\item process input/output declarations inconsistent with graph flows or ports;
\item disconnected data stores;
\item decomposition references pointing to missing graph layers;
\item cycles, reported separately from path generation.
\end{itemize}

\subsection{System Functions}

The implemented system provides these user-facing functions:

\begin{itemize}
\item import standard JSON and SOFL `.cdfd' files;
\item infer or accept explicit start and end nodes;
\item generate atomic source-to-sink paths;
\item generate structured concurrent paths;
\item analyze path relations such as parallel, exclusive, and joined-output;
\item expand multi-level process decompositions for atomic paths;
\item render SOFL-style SVG diagrams;
\item provide Web UI, Web API, and CLI output;
\item export text, JSON, CSV, and Markdown results.
\end{itemize}

\section{Key Technical Details}

This section explains the model design and the algorithms used for different CDFD cases.

\subsection{Model Design}

The common model separates graph facts from algorithm results. `Node' and `Edge' store the imported CDFD. `ProcessSpec' and `PortSpec' store process interfaces. `GraphStructure' stores explicit CDFD structures such as parallel and choice. `PathResult' stores one atomic path, including nodes, edge ids, edge sources, edge targets, edge data, output data, preconditions, and control conditions. `ConcurrentPathResult' stores a tree with `node', `sequential', and `parallel' elements plus a flat summary for export.

This design is intentionally conservative. If a process has multiple output ports and no explicit parallel/fork structure, the algorithm treats different ports as alternatives. If several outgoing edges are assigned to the same output port, they can be produced together by one selected port.

\subsection{Atomic Path Algorithm}

The atomic path generator performs depth-first search over traversable data-flow edges. At each step it updates available data, activated nodes, and collected conditions. A candidate edge is accepted only if the target process can be activated under the current data and port state.

The generator supports two cycle strategies. In the default `simple' strategy, a node cannot be repeated in the same path, so termination follows from the finite node set. In the `max-depth' strategy, repeated nodes are allowed up to a user-provided depth bound. A global `max\_paths' limit prevents accidental combinational explosion.

\subsection{Parallel Forks}

When a node has several outgoing branches, the algorithm must decide whether they are alternatives or concurrent branches. Explicit `parallel', `broadcast', `fork', or `separate' structures mark parallel branches. Explicit `choice', `condition', `select', or `non-determinism' structures mark alternatives. Without explicit structure, outgoing branches from different output ports are treated as alternatives. Outgoing branches from the same selected output port may be grouped as a concurrent branch if their conditions do not conflict.

\subsection{Join and Synchronized Multi-Input}

Join behavior is handled by token search. A token may reach a process before all required input data is available. In that case, the algorithm records the data that has arrived but delays activation of the join target. The target becomes active only when its required input port is ready. This supports examples where two branches must both finish before a downstream process can run.

For synchronized multi-start graphs, the initial state can contain multiple tokens, such as one external input token and one data-store token. The resulting concurrent notation can be printed as `[IN || STORE] -> Process -> OUT'.

\subsection{Mutually Exclusive Ports}

Multiple input ports are treated as alternatives. For example, a login process may accept either a password port \(\{userAccount, passWord\}\) or a token port \(\{token\}\). The password port cannot fire when only `passWord' is available, because the port requires both data items. The token port can fire independently. This is why a single process node may still produce different legal paths depending on which port is activated.

Multiple output ports are also alternatives by default. This prevents the tool from incorrectly treating all outgoing edges of a process as simultaneous outputs. Parallel output must be represented either by a shared selected output port or by an explicit CDFD parallel/fork structure.

\subsection{Multi-Level Process Decomposition}

A process specification may include a `decom' field that points to a lower-level graph. During atomic path generation, the multilevel module replaces that process with the paths of its decomposition graph and connects the parent graph edges to the child graph inputs and outputs. The current implementation supports stable atomic expansion across several graph layers. Fully unified concurrent expansion across nested decomposition graphs remains an open limitation.

\subsection{Visualization and Output}

The SVG renderer uses SOFL-like process frames, data-store boxes, condition diamonds, solid data-flow arrows, dashed control-flow arrows, and imported layout coordinates when available. JSON graphs without layout data are automatically arranged. The Web UI focuses on the imported graph, generated paths, concurrent paths, path relations, and consistency warnings. The CLI exposes the same analysis for reproducible command-line use.

\section{Case Studies}

Table~\ref{tab:cases} summarizes representative examples used in testing and demonstration.

\begin{table*}[t]
\centering
\caption{Representative CDFD Path Generation Cases}
\label{tab:cases}
\begin{tabular}{|p{0.15\textwidth}|p{0.18\textwidth}|p{0.27\textwidth}|p{0.30\textwidth}|}
\hline
Input & Main CDFD Feature & Generated Result & Meaning \\
\hline
`cdfd\_v1.json' & Fork after process A & Two atomic paths; B and C are reported as parallel branches. Current notation includes `IN -> A -> [B || C] -> OUT\_X4 -> OUT\_X5'. & The tool recognizes independent branches after A. The terminal-output notation is still conservative when several sink nodes appear. \\
\hline
`join.json' & Fork followed by join & Two atomic paths and one concurrent path: `IN -> Split -> [L || R] -> Combine -> OUT'. & Combine is activated only after both `l\_done' and `r\_done' become available. \\
\hline
`data\_store.json' & Synchronized multi-input & Concurrent notation: `[IN || PROFILE\_STORE] -> BuildResponse -> OUT'. & External request data and stored profile data are both needed before the response process runs. \\
\hline
`multilevel.json' & Process decomposition & Three expanded atomic paths across four graph layers. & A1, A3, and A33 are replaced by their lower-level CDFDs; state/control nodes s1 and s2 become path conditions. \\
\hline
`port\_alternatives.json' & Mutually exclusive input ports & Two legal paths: one through the token port and one through the password port. & The password port requires both `userAccount' and `passWord'; `passWord' alone cannot activate Login. \\
\hline
`xuexitong.cdfd' & Real SOFL `.cdfd' import & Four atomic paths and an exclusive path relation. & The importer preserves process, data store, condition, control-flow, and layout information from the SOFL drawing. \\
\hline
`loop.json' & Cyclic CDFD structure & Simple mode avoids repeated nodes; max-depth mode allows bounded loop exploration. & Termination is controlled by path strategy and depth/path limits. \\
\hline
\end{tabular}
\end{table*}

The following output from `join.json' illustrates the intended concurrent path notation:
\begin{equation}
  IN \rightarrow Split \rightarrow [L \parallel R] \rightarrow Combine \rightarrow OUT.
\end{equation}
This single structured path corresponds to two legal atomic projections, one through \(L\) and one through \(R\). The join process is not activated by either projection alone; it requires both branch outputs.

\section{Uncertainties and Limitations}

Two points are not fully settled and should be stated explicitly.

First, SOFL `.cdfd' files store connector and coordinate information in tool-specific ways. The importer preserves connector indexes and infers missing endpoints, but more real SOFL files are needed to verify every port-mapping case.

Second, multi-level atomic expansion is implemented, but fully unified concurrent expansion across nested decomposition graphs is not complete. A future version should replace a decomposed process with a concurrent subtree when the child graph contains parallel structure.

\section{Conclusion}

This paper presented a CDFD path generation method and tool that covers the required tasks: defining a complete input format, parsing CDFD data into a common graph model, checking legality and consistency, generating atomic and concurrent paths, handling process ports and different algorithmic cases, and visualizing results. The main idea is to treat CDFD path generation as a semantic graph problem. Atomic paths capture source-to-sink data-flow traces, while concurrent paths capture independent branches and synchronized joins. The port-aware rules prevent the common error of treating a multi-port process as if all of its interfaces were active at once. The implementation and case studies show that the approach can support JSON inputs, SOFL `.cdfd' imports, multilevel examples, port alternatives, joins, and cyclic graphs.

\bibliographystyle{template/IEEEtran}
\bibliography{references}

\end{document}
